\chapter{緒論}\label{chap:introduction}
\begingroup
\setlength{\emergencystretch}{2em}
\clubpenalty=10000
\widowpenalty=10000

\section{研究背景}

分類模型常以過去累積的資料訓練，再用來判斷較晚時間出現的樣本。這個做法有一項前提：過去學到的關係，在後續資料中仍具有參考價值。然而，資料來源、觀測對象及外在環境都可能改變，舊模型的表現也可能隨之變動。概念漂移（concept drift）研究便關心輸入與目標之間的關係如何隨時間改變，以及模型應如何因應這些變化\autocite{gama2014drift}。對具有時間資訊的分類問題而言，訓練資料要保留多久、近期資料要如何使用，因而成為模型建立時需要處理的問題。

另一項挑戰是類別不平衡。當研究關注的正類只占少數，多數類所提供的訊息可能主導模型的學習。即使整體正確率較高，也未必表示模型能辨認研究關心的樣本。例如，正類占比很低時，將大部分樣本判為負類便可能取得看似不錯的正確率。Saito 與 Rehmsmeier\autocite{saito2015pr}指出，不平衡分類的評估需要留意正類預測的可靠程度，不能只從 ROC 圖的外觀判斷模型是否適合使用。

時間與不平衡問題也可能同時出現。較長的歷史期間通常提供較多訓練紀錄，卻不一定與近期資料具有相同特性；近期資料較接近測試期間，但少數類樣本可能不足。若直接合併全部資料，較多的歷史樣本可能影響模型的學習方向；若完全捨棄歷史資料，又可能失去仍有參考價值的分類資訊。本研究因此從歷史與近期資料的運用方式出發，比較分別建模、合併重訓及模型集成的表現。

本研究以公開的美國企業年度財務資料作為實證案例。企業財務資訊具有明確年度，且資料中的破產類別相對較少，適合用來觀察不同訓練期間與不平衡處理方式的差異。資料的研究背景來自 Lombardo 等人\autocite{lombardo2022bankruptcy}，實際使用作者公開更新的版本\autocite{sowide-bankruptcy-data}。本地資料共含 1999--2018 年的 78,682 筆公司年度紀錄、8,971 個匿名公司識別碼與 18 個財務特徵，其中正類紀錄占 6.63\%。資料版本、欄位及標籤的核對結果於第三章說明。

需要先界定的是，公開資料的公司狀態在同一公司所有年度保持一致，並未附上可逐筆核對的破產日期。因此，本研究沿用公司最終存續或破產狀態作為標籤，以較早年度資料訓練、較晚年度資料測試，研究對象是既定時間分區下的公司狀態分類。本文使用「破產預測」指稱此應用背景，實證結果不解釋為提前一年預測破產事件的能力。

\section{研究動機}

改善少數類辨識能力，可以從資料採樣、模型組成及特徵空間等方向著手。採樣改變模型接觸到的訓練資料；集成結合不同模型的判斷；特徵選取則調整模型能使用的資訊。這些方法處理的問題不同，搭配使用時的效果也需要透過對照實驗確認，不能只因某一方法在其他研究有效，就預期相同結論會出現在本資料中。

本研究首先關心資料期間的取捨。若近期模型的分數較高，可能與資料較接近測試期間有關，也可能同時受到樣本數、採樣與驗證設定的影響。因而，除了比較歷史資料模型 Old、近期資料模型 New 與合併資料重訓模型 Retrain，也有必要記錄各種分類器與採樣設定，避免將單一分數差異簡化為時間因素造成的效果。

其次，歷史與近期模型是否值得一起保留，不能只看各自的分數。若模型在不同樣本上各有所長，集成可能提供互補資訊；若模型的錯誤相近，加入更多模型也不一定有幫助。靜態集成在推論前固定成員，動態選擇則依待判斷樣本的鄰近資料估計模型能力。兩種做法都具有合理出發點，實際表現仍須在相同模型池與驗證規則下比較。本研究以固定分類器建立 Old／New 六模型池，使比較能集中於模型組成與選擇方式。

最後，刪減特徵可能移除重複或干擾資訊，也可能刪掉少數類辨識所需的線索。動態方法另外依賴特徵空間中的鄰近關係，因此特徵選取除了影響基礎分類器，也可能改變其局部能力估計。本研究將特徵選取安排為另一項實驗，沿用既有集成條件，分別比較不同方法及保留數，而不預設保留較少特徵必然較好。

\section{研究目的與研究問題}

本研究的目的，是建立具有時間資訊之不平衡分類比較流程，說明資料期間、集成方式及特徵選取如何影響分類結果。流程以通用的樣本、時間、正類及特徵表示；企業資料提供本次實證案例。具體研究問題如下。

\begin{enumerate}
\item \textbf{研究問題一：不同訓練期間與採樣方式下，單一分類器呈現哪些差異？}以已保存的 Baseline 實驗比較 Old、New 與 Retrain，整理不同分類器及採樣的 AUC、F1-score 與 Recall。此部分提供描述性基準，並明列歷史調參條件的差異。
\item \textbf{研究問題二：固定 Old／New 模型池後，靜態與動態集成的表現是否不同？}使用相同分類器、六個基礎模型及固定驗證資料，比較單一期間集成、六模型等權集成、驗證選出的靜態集成，以及動態分類器與動態集成選擇。
\item \textbf{研究問題三：在相同集成流程中加入特徵選取，是否改變區辨能力與少數類辨識表現？}比較四種特徵選取方法及兩種保留數，並與同方法、同種子、同測試樣本的未選取特徵結果配對，檢視效能與特徵精簡程度的關係。
\end{enumerate}

上述問題安排為兩項實驗。實驗一包含 Baseline 與固定年份集成比較，實驗二則為特徵選取後的集成比較。AUC 用來觀察整體分數排序，F1-score 與 Recall 用來補充正類分類結果；靜態成員與分類門檻均由驗證資料決定，測試資料不參與選擇。第四章依此順序呈現結果，並區分描述性差異與經多重比較校正後的統計證據。

\section{研究範圍與貢獻}

本研究的貢獻在於整理可追溯的比較流程與實證證據，並非提出新的分類演算法。首先，將資料期間與採樣來源納入 Old／New 模型池的組成，使不同集成方法可以使用相同成員進行比較。其次，在固定年份集成設定中分開模型擬合、動態選擇參考及候選選擇用途，讓模型成員與門檻的決策來源清楚可查。再者，特徵選取在採樣前使用允許的擬合資料執行，並以完整特徵結果作為配對對照，觀察刪減資訊後的變化。

結果同時列出合併四年、各年度與十個種子的分數，並以公司為重抽樣單位比較差值及校正多重檢定。同公司的年度紀錄、隨機種子與重疊年份邊界不視為獨立樣本，推論範圍亦限於既定資料與模型。

本次實證限於一份公開資料。同一公司可能跨訓練與測試年度出現，因此評估不屬於完全未見公司的驗證。主要集成與 FS 採固定年份切分，測試期間不更新模型；第三章另以滑動視窗圖說明一般時間配置概念。各比較所需的控制條件與推論限制於第三、四章交代。

\section{論文章節安排}

本文共五章。第一章交代背景、研究問題與範圍；第二章整理相關文獻；第三章說明通用架構、資料、兩項實驗與驗證程序；第四章呈現結果與討論；第五章整理結論與後續方向。補充設定、完整結果與檢定細節置於附錄。
\endgroup
